%% sections/03-pattern.tex

\section{The Arc-Completion Pattern}
\label{sec:pattern}

\subsection{Design Template}

Each NPC designed with arc-completion potential receives an arc-completion section
in their NPC file with four fields:

\begin{description}
  \item[Arc-completion text:] The exact line of dialogue, act, or gesture the NPC
    delivers. Specific, brief, and not expandable by the DM without compromising the
    moment. Example: \textit{``Veth: `I don't know if I am doing it right. I do it
    anyway.' --- said once, then back to work.''}

  \item[Firing condition:] The precise party behavior that triggers the arc-completion.
    Not ``if the party earns Veth's trust'' but ``if the party asks Veth what he is
    trying to understand.'' The condition is specific enough that the DM (or AI system)
    knows with certainty when it has been met.

  \item[DM discipline note:] Explicit instruction on what NOT to do. Typically: do not
    extend the moment, do not have the NPC expand on the arc-completion line, do not
    invite the party to respond in a way that disrupts the landing.

  \item[Non-firing note (optional):] The NPC's behavior if the firing condition is NOT
    met. Often: the NPC remains at their current emotional state; the arc-completion is
    simply absent, not replaced by a lesser version.
\end{description}

Table~\ref{tab:examples} shows three representative arc-completion designs from the
corpus.

\begin{table*}[t]
\caption{Three Representative NPC Arc-Completion Designs}
\label{tab:examples}
\begin{tabular}{@{}lp{4cm}p{4cm}p{3cm}@{}}
\toprule
NPC & Arc-Completion Text & Firing Condition & Campaign/Session \\
\midrule
Mig & She steps through, and does not look back. [No line. Act only.] & Party opens the inner door cooperatively, making no comment on Mig's hesitation. & C3-S1 \\
Gorret & ``At least you say what you are.'' --- said once; then stepped aside. & A party member gives an honest answer to Gorret's direct question ``why do you need to go up?'' & C3-S2 \\
Saren Coldforge & [Takes Asholt's token from his hand across the parley table. Stands. Says nothing.] & Asholt extends the token with his right hand, and Saren is present at the table when he does. & C4-S7 \\
\bottomrule
\end{tabular}
\end{table*}

\subsection{Character Logic vs.\ DM Scripting}

The distinction between \textit{character logic} and \textit{DM scripting} is central
to the pattern's value. DM scripting occurs when the DM decides independently that
``now is the time'' for an NPC moment and delivers it regardless of party behavior.
Character logic occurs when the NPC's arc-completion fires because the party's behavior
met the pre-specified condition.

In practice, the distinction is observable in the session log: character-logic firings
are preceded by specific party actions that match the firing condition; DM-scripted
firings are not. In the Marathon corpus, all 40+ logged instances were character-logic
firings. This is partly a design artifact (the AI system does not fire arc-completions
without checking firing conditions) but also reflects that the firing conditions were
well-specified enough to be checkable.

\subsection{The Small-Moment Principle}

A consistent finding across the corpus is that arc-completions are smaller than
expected. Not a speech, not a monologue, but one line, one act, one gesture. This is
a design requirement, not a default: the arc-completion text is constrained to be
minimal by the DM discipline note (``do not extend the moment'').

Table~\ref{tab:small-moments} shows the word counts of the top-10 most impactful
arc-completion lines in the corpus (measured by panel AL citation frequency).

\begin{table}[t]
\caption{Word Counts of Top-10 Most Panel-Cited Arc-Completion Moments}
\label{tab:small-moments}
\begin{tabular}{@{}llrl@{}}
\toprule
NPC & Words & Campaign & What makes it work \\
\midrule
Mig & 0 (act only) & C3-S1 & Physical threshold crossing; no words needed \\
Veth & 14 & C3-S2 & Names the doubt and the practice together \\
Gorret & 9 & C3-S2 & Acknowledges without approving \\
Pella & 2 & C1-S2 & Campaign-long anticipation in 2 words \\
Saren (token) & 0 (act only) & C4-S7 & 7 sessions of investment released in a gesture \\
Asholt & 12 & C4-S7 & Names what was wrong and what was right in same sentence \\
Bessa Tor & 18 & C5-S1 & Gives key without conditions; only explanation: ``You asked right'' \\
Lenne (suspension) & 0 (act: didn't write) & C2-S3 & First documentation-suspension in C2 \\
Mira S (naming) & 21 & C1-S3 & 30 years of refusal ended in writing a name \\
Velantha (journal) & 25 & C3-S7 & Entire arc completion in one final journal entry \\
\bottomrule
\end{tabular}
\end{table}

The median word count for the top-10 arc-completions is 7 words. Two are wordless acts.
None exceeds 25 words. The session-memory-worthy moments in the corpus are
disproportionately the smallest.

This is counterintuitive from a narrative design perspective: the conventional assumption
is that important moments deserve more space. The arc-completion pattern inverts this:
important moments deserve \textit{less} space, because the space has already been
accumulated by the preceding sessions. The moment is the hinge, not the speech.

\subsection{NPC Design Discipline Requirements}

The pattern requires three DM discipline commitments that are frequently violated
in practice:

\paragraph{Do not deliver early.}
The most common discipline failure is delivering the arc-completion before the firing
condition is met. A DM who can see that ``the party is getting close'' to the condition
may deliver a partial arc-completion, diluting the eventual full firing. The discipline
note must state: ``Do not deliver any version of this line before the exact firing
condition is met.''

\paragraph{Do not extend.}
The second most common discipline failure is adding NPC dialogue after the arc-completion
to explain or contextualize it. The arc-completion text is the completion; the NPC
does not explain it. The DM discipline note must state: ``After delivering the arc-completion
text, the NPC is done. Do not continue speaking for the NPC.''

\paragraph{Do not soften.}
Some NPC arc-completions are uncomfortable: a morally grey character acknowledges
something without becoming sympathetic; an antagonist's closing line is not a conversion.
The discipline note must state: ``Do not soften the arc-completion to make the NPC
more sympathetic. The line is what it is. Deliver it as written.''

In the AI-simulated corpus, all three discipline commitments were enforced
mechanically by the session runner, which reads the NPC file's arc-completion
section and fires it only when the firing condition is confirmed. Human DMs must
enforce these as volitional commitments.

The small-moment principle reflects a theory of emotional payoff in narrative: the
most emotionally resonant moments are often the most compressed. When the party has
been working toward a moment across multiple scenes, the moment itself needs very
little space --- the accumulated context does the work. Extending the moment dissipates
this accumulated energy rather than releasing it.

Empirically, the smallest arc-completions in the corpus scored highest on Atmospheric
Landing. Mig's wordless threshold-crossing (C3-S1) and Saren's silent token-acceptance
(C4-S7) were both noted in panel reviews as session-memory-worthy moments. Extended
NPC responses in the same sessions were not noted.
